\documentclass[11pt,aspectratio=169]{beamer}

\usepackage[utf8]{inputenc}
\usepackage[T5]{fontenc}
\usepackage[vietnamese]{babel}
\shorthandoff{"}
\AtBeginDocument{\shorthandoff{"}}
\pdfstringdefDisableCommands{\def\\{ }}

\usepackage{graphicx}
\usepackage{amsmath,amssymb,amsfonts,mathtools,bm}
\usepackage{booktabs,tabularx,array,multirow}
\usepackage{tikz}
\usetikzlibrary{arrows.meta,positioning,shapes.geometric,calc,fit,backgrounds}
% Giữ thuật ngữ chuyên môn bằng tiếng Anh; phần giải thích dùng tiếng Việt.
% Ngôn ngữ thiết kế theo PDF Pretraining: Madrid, 16:9, Latin Modern,
% xanh tím #3333B2, cam #C46B22, nền trắng và các khối tím nhạt.
\usetheme{Madrid}
\usefonttheme{serif}
\usepackage{lmodern}

\definecolor{refblue}{HTML}{3333B2}
\definecolor{reforange}{HTML}{C46B22}
\definecolor{refgray}{HTML}{727272}
\definecolor{refsoft}{HTML}{E9E9F3}
\colorlet{navyblue}{refblue}
\colorlet{stanfordred}{reforange}
\colorlet{tealgreen}{refblue}
\colorlet{amberorange}{reforange}
\colorlet{slate}{refgray}
\colorlet{softblue}{refsoft}
\colorlet{softgreen}{refsoft}
\colorlet{softred}{reforange!7!white}
\colorlet{softyellow}{reforange!7!white}

\setbeamercolor{normal text}{fg=black,bg=white}
\setbeamercolor{structure}{fg=refblue}
\setbeamercolor{title}{fg=white,bg=refblue}
\setbeamercolor{frametitle}{fg=white,bg=refblue}
\setbeamercolor{block title}{fg=white,bg=refblue!75!black}
\setbeamercolor{block body}{fg=black,bg=refsoft}
\setbeamercolor{block title example}{fg=white,bg=refblue!75!black}
\setbeamercolor{block body example}{fg=black,bg=refsoft}
\setbeamercolor{block title alerted}{fg=white,bg=refblue!75!black}
\setbeamercolor{block body alerted}{fg=black,bg=refsoft}
\setbeamercolor{author in head/foot}{fg=white,bg=refblue!50!black}
\setbeamercolor{title in head/foot}{fg=white,bg=refblue!75!black}
\setbeamercolor{date in head/foot}{fg=white,bg=refblue}
\setbeamerfont{frametitle}{size=\large,series=\mdseries}
\setbeamerfont{block title}{size=\normalsize,series=\mdseries}
\setbeamerfont{title}{size=\Large,series=\mdseries}
\setbeamerfont{subtitle}{size=\small}
\setbeamerfont{author}{size=\small}
\setbeamerfont{institute}{size=\footnotesize}
\setbeamerfont{date}{size=\footnotesize}
\setbeamerfont{author in head/foot}{size=\tiny}
\setbeamerfont{title in head/foot}{size=\tiny}
\setbeamerfont{date in head/foot}{size=\tiny}
\setbeamertemplate{navigation symbols}{}
\setbeamertemplate{caption}[numbered]
\setbeamersize{text margin left=0.8cm,text margin right=0.8cm}

\newsavebox{\fitlinebox}
\newcommand{\fitline}[2]{%
  \sbox{\fitlinebox}{#2}%
  \ifdim\wd\fitlinebox>#1\relax
    \resizebox{#1}{!}{\usebox{\fitlinebox}}%
  \else\usebox{\fitlinebox}\fi
}
\newcommand{\fitmath}[1]{\fitline{\linewidth}{$\displaystyle #1$}}

\setbeamertemplate{footline}{%
  \leavevmode\hbox{%
    \begin{beamercolorbox}[wd=.333333\paperwidth,ht=2.25ex,dp=1ex,center]{author in head/foot}%
      \usebeamerfont{author in head/foot}Post-training (2020--2026)
    \end{beamercolorbox}%
    \begin{beamercolorbox}[wd=.333333\paperwidth,ht=2.25ex,dp=1ex,center]{title in head/foot}%
      \usebeamerfont{title in head/foot}Post-training of Language Models
    \end{beamercolorbox}%
    \begin{beamercolorbox}[wd=.333334\paperwidth,ht=2.25ex,dp=1ex,right]{date in head/foot}%
      \usebeamerfont{date in head/foot}\insertframenumber{} / \inserttotalframenumber\hspace*{2ex}
    \end{beamercolorbox}%
  }%
}

% Tiêu đề phụ phẳng và câu dẫn, theo cách trình bày trong tài liệu mẫu.
\newcommand{\subhead}[1]{%
  \par\vspace{0.10cm}{\small\bfseries\color{refblue}#1\par}\vspace{0.06cm}%
}
\newcommand{\limithead}[1]{\subhead{#1}}
\newcommand{\motivehead}[1]{\subhead{#1}}
\newcommand{\lead}[1]{%
  {\small #1\par}\vspace{0.16cm}%
}
\newcommand{\source}[1]{%
  \par\vspace{0.12cm}{\scriptsize\color{refgray}#1\par}%
}
\newcommand{\takeaway}[1]{%
  \begin{beamercolorbox}[rounded=true,shadow=true,sep=4pt]{block body}
    \small #1
  \end{beamercolorbox}%
}

\AtBeginDocument{%
  \setlength{\abovedisplayskip}{5pt}%
  \setlength{\belowdisplayskip}{5pt}%
  \setlength{\abovedisplayshortskip}{3pt}%
  \setlength{\belowdisplayshortskip}{3pt}%
}

% Trang chuyển phần: nhãn cam, tiêu đề xanh và năm đoạn tiến trình.
\newcommand{\sectionintro}[4]{%
  \section[#2]{#2}
  \begin{frame}[plain,noframenumbering,t]
    \vspace{1.0cm}
    {\large\bfseries\color{reforange}PHẦN #1 / 04\par}
    \vspace{0.55cm}
    {\fontsize{28}{32}\selectfont\bfseries\color{refblue}#2\par}
    \vspace{0.20cm}
    {\normalsize\color{refgray}#3\par}
    \vspace{0.65cm}
    {\small\color{refgray}\textbf{Nội dung:} #4\par}
    \vspace{0.40cm}
    \foreach \step in {1,...,4}{%
      \ifnum\step=#1\relax\color{reforange}\else\color{refblue!20}\fi
      \rule{2.30cm}{2.2pt}\hspace{0.10cm}%
    }\par
    \vspace{0.15cm}
    {\scriptsize\color{refgray}Bốn giai đoạn của Post-training: từ SFT đến GRPO / RLVR\par}
  \end{frame}%
}

% Sơ đồ giữ nguyên quan hệ và nhãn; đồng bộ bảng màu và chữ serif.
\tikzset{
  >=Stealth,
  data/.style={rectangle,rounded corners=2pt,draw=reforange!70,line width=0.65pt,
    fill=softyellow,text width=3.0cm,align=center,minimum height=0.95cm,
    inner sep=4pt,font=\scriptsize\rmfamily},
  trainmodel/.style={rectangle,rounded corners=2pt,draw=refblue!70,line width=0.8pt,
    fill=refsoft,text width=3.1cm,align=center,minimum height=1.05cm,
    inner sep=4pt,font=\scriptsize\rmfamily\bfseries},
  frozenmodel/.style={rectangle,rounded corners=2pt,draw=refgray,line width=0.65pt,dashed,
    fill=gray!7,text width=3.1cm,align=center,minimum height=1.05cm,
    inner sep=4pt,font=\scriptsize\rmfamily},
  rewardbox/.style={rectangle,rounded corners=2pt,draw=refblue!60,line width=0.65pt,
    fill=refsoft,text width=3.0cm,align=center,minimum height=0.95cm,
    inner sep=4pt,font=\scriptsize\rmfamily},
  lossbox/.style={rectangle,rounded corners=2pt,draw=reforange!80,line width=0.8pt,
    fill=softyellow,text width=3.2cm,align=center,minimum height=1.0cm,
    inner sep=4pt,font=\scriptsize\rmfamily\bfseries},
  flow/.style={->,line width=0.75pt,draw=refblue!75,rounded corners=2pt},
  gradflow/.style={->,line width=0.8pt,draw=reforange,dashed,rounded corners=2pt}
}


\DeclareMathOperator{\sen}{sen}
\DeclareMathOperator{\tg}{tg}
\DeclareMathOperator*{\argmax}{arg\,max}
\DeclareMathOperator*{\argmin}{arg\,min}
\DeclareMathOperator{\E}{\mathbb{E}}
\DeclareMathOperator{\DKL}{D_{\mathrm{KL}}}
\DeclareMathOperator{\clip}{\mathrm{clip}}
\DeclareMathOperator{\mean}{\mathrm{mean}}
\DeclareMathOperator{\std}{\mathrm{std}}
\newcommand{\cancel}[1]{\textcolor{reforange}{\underbrace{#1}_{\text{Triệt tiêu }=0}}}

\title[Post-training of Language Models]{Post-training\\of Language Models}
\author[Nhóm thực hiện]{%
  Hà Thanh Phong -- 24520024\\
  Hà Xuân Thiện -- 24520031\\
  Trần Quang Trường -- 24521901}
\institute{University of Information Technology -- VNUHCM}
\date{\today}
\bibliographystyle{apalike}

\begin{document}

% Cách 1: Đánh số trong tiêu đề block: "Công thức (1): ..."
\begin{frame}{Cách 1: Đánh số trong tiêu đề Block}
  \setlength{\abovedisplayskip}{2pt}%
  \setlength{\belowdisplayskip}{2pt}%
  \lead{Tối đa hóa phần thưởng kết hợp KL regularization (Ouyang et al., 2022).}
  \vspace{-0.15cm}
  \begin{block}{Mục tiêu RLHF có điều chuẩn KL \hfill (1)}
    \[
      \fitmath{%
      \max_\theta \; J_{\text{RLHF}}(\theta) = \E_{x \sim \mathcal{D},\; y \sim \pi_\theta} \big[ R_{\text{total}}(x, y) \big], \quad \text{với } R_{\text{total}}(x, y) = r_\phi(x, y) - \beta \log \frac{\pi_\theta(y \mid x)}{\pi_{\text{ref}}(y \mid x)}
      }
    \]
  \end{block}
  \vspace{-0.15cm}
  \begin{block}{Policy Gradient với hàm Advantage $A(x, y)$ \hfill (2)}
    \[
      \fitmath{%
      \nabla_\theta J_{\text{RLHF}}(\theta) = \E_{x \sim \mathcal{D},\; y \sim \pi_\theta} \Big[ \underbrace{\big(R_{\text{total}}(x, y) - b(x)\big)}_{\text{Advantage } A(x, y)} \, \nabla_\theta \log \pi_\theta(y \mid x) \Big]
      }
    \]
  \end{block}
  \vspace{-0.10cm}
  \subhead{Ý nghĩa của Baseline $b(x)$ và Advantage $A(x, y)$}
  \footnotesize
  \begin{itemize}\setlength{\itemsep}{2pt}
    \item \textbf{Baseline $b(x)$:} Mốc kỳ vọng của prompt $x$ (giúp giảm phương sai gradient mà không gây chệch).
    \item \textbf{Advantage $A(x, y) = R_{\text{total}} - b(x)$:} Đo mức vượt trội so với mốc $b(x)$ ($A > 0$: tăng $\uparrow$; $A < 0$: giảm $\downarrow$).
  \end{itemize}
  \par\vspace{0pt}
\end{frame}

% Cách 2: Đánh số bên phải công thức
\begin{frame}{Cách 2: Đánh số bên phải trong công thức}
  \setlength{\abovedisplayskip}{2pt}%
  \setlength{\belowdisplayskip}{2pt}%
  \lead{Tối đa hóa phần thưởng kết hợp KL regularization (Ouyang et al., 2022).}
  \vspace{-0.15cm}
  \begin{block}{Mục tiêu RLHF có điều chuẩn KL}
    \[
      \fitmath{%
      \max_\theta \; J_{\text{RLHF}}(\theta) = \E_{x \sim \mathcal{D},\; y \sim \pi_\theta} \big[ R_{\text{total}}(x, y) \big], \quad \text{với } R_{\text{total}}(x, y) = r_\phi(x, y) - \beta \log \frac{\pi_\theta(y \mid x)}{\pi_{\text{ref}}(y \mid x)} \qquad (1)
      }
    \]
  \end{block}
  \vspace{-0.15cm}
  \begin{block}{Policy Gradient với hàm Advantage $A(x, y)$}
    \[
      \fitmath{%
      \nabla_\theta J_{\text{RLHF}}(\theta) = \E_{x \sim \mathcal{D},\; y \sim \pi_\theta} \Big[ \underbrace{\big(R_{\text{total}}(x, y) - b(x)\big)}_{\text{Advantage } A(x, y)} \, \nabla_\theta \log \pi_\theta(y \mid x) \Big] \qquad (2)
      }
    \]
  \end{block}
  \vspace{-0.10cm}
  \subhead{Ý nghĩa của Baseline $b(x)$ và Advantage $A(x, y)$}
  \footnotesize
  \begin{itemize}\setlength{\itemsep}{2pt}
    \item \textbf{Baseline $b(x)$:} Mốc kỳ vọng của prompt $x$ (giúp giảm phương sai gradient mà không gây chệch).
    \item \textbf{Advantage $A(x, y) = R_{\text{total}} - b(x)$:} Đo mức vượt trội so với mốc $b(x)$ ($A > 0$: tăng $\uparrow$; $A < 0$: giảm $\downarrow$).
  \end{itemize}
  \par\vspace{0pt}
\end{frame}

% Cách 3: Ghi rõ "Công thức (1): ..." trong block title
\begin{frame}{Cách 3: Ghi tên công thức}
  \setlength{\abovedisplayskip}{2pt}%
  \setlength{\belowdisplayskip}{2pt}%
  \lead{Tối đa hóa phần thưởng kết hợp KL regularization (Ouyang et al., 2022).}
  \vspace{-0.15cm}
  \begin{block}{(1) Mục tiêu RLHF có điều chuẩn KL}
    \[
      \fitmath{%
      \max_\theta \; J_{\text{RLHF}}(\theta) = \E_{x \sim \mathcal{D},\; y \sim \pi_\theta} \big[ R_{\text{total}}(x, y) \big], \quad \text{với } R_{\text{total}}(x, y) = r_\phi(x, y) - \beta \log \frac{\pi_\theta(y \mid x)}{\pi_{\text{ref}}(y \mid x)}
      }
    \]
  \end{block}
  \vspace{-0.15cm}
  \begin{block}{(2) Policy Gradient với hàm Advantage $A(x, y)$}
    \[
      \fitmath{%
      \nabla_\theta J_{\text{RLHF}}(\theta) = \E_{x \sim \mathcal{D},\; y \sim \pi_\theta} \Big[ \underbrace{\big(R_{\text{total}}(x, y) - b(x)\big)}_{\text{Advantage } A(x, y)} \, \nabla_\theta \log \pi_\theta(y \mid x) \Big]
      }
    \]
  \end{block}
  \vspace{-0.10cm}
  \subhead{Ý nghĩa của Baseline $b(x)$ và Advantage $A(x, y)$}
  \footnotesize
  \begin{itemize}\setlength{\itemsep}{2pt}
    \item \textbf{Baseline $b(x)$:} Mốc kỳ vọng của prompt $x$ (giúp giảm phương sai gradient mà không gây chệch).
    \item \textbf{Advantage $A(x, y) = R_{\text{total}} - b(x)$:} Đo mức vượt trội so với mốc $b(x)$ ($A > 0$: tăng $\uparrow$; $A < 0$: giảm $\downarrow$).
  \end{itemize}
  \par\vspace{0pt}
\end{frame}

\end{document}
